% NOTATION -- master table, complete. The functional-analysis and spectral-
% truncation symbols of Part I, among them the retained mode set Z_{k_max} and
% the smoothing order alpha of Ch. 5, whose Appendix B diffusivity namesake is
% disambiguated in the same row; the tensor-layout symbols of Part II, used by
% Chapters 10--16; and the physical symbols of Part III, used by
% Chapters 17--24. Every symbol whose name is reused across parts carries its
% ambiguity in its own row, as the chapters themselves print it.
\chapter*{Notation and Conventions}
\addcontentsline{toc}{chapter}{Notation and Conventions}
\markboth{Notation and Conventions}{Notation and Conventions}

Symbols are used consistently throughout the book. The conventions below are
established in Part~I and reused without restatement in Parts~II and~III, where
the second block of the table records the array conventions every implementation
is written in and the third the physical quantities the application chapters
solve for. Where the same letter carries two meanings in different parts, the row
for the later meaning names the earlier one where space allows --- the reuses the
reader is most likely to meet are $p$ (integrability exponent, DeepONet latent
width, parameter count), $\sigma$ (activation, pointwise spread, Sobolev order in
$H^\sigma$), $m$
(sensor count, modes per band edge), $P$ (the DeepONet query-point count, and the projectors $P_N$,
$P_{k_{\max}}$), and $L$ (layer count, body length) --- but the flagging is a
promise about visibility, not a proof that no symbol anywhere does double duty.

% {\raggedright\arraybackslash} on the definition column: the cells are narrow
% single paragraphs and justification leaves rivers and underfull lines in them
% (pattern used for every other table column in the book).
{\small
\begin{xltabular}{\textwidth}{@{}l>{\raggedright\arraybackslash}X@{}}
\caption{\textbf{Notational conventions: the infinite-dimensional objects, the tensors that realize them, and the physical symbols the applications are written in.}}
\label{tab:notation-master}\\
\toprule
\textbf{Symbol} & \textbf{Mathematical Definition and Context} \\
\midrule
\endfirsthead
\multicolumn{2}{r}{\textit{Table~\ref{tab:notation-master}, continued}}\\
\toprule
\textbf{Symbol} & \textbf{Mathematical Definition and Context} \\
\midrule
\endhead
\midrule
\multicolumn{2}{r}{\textit{Table~\ref{tab:notation-master}, continued on the next page}}\\
\endfoot
\bottomrule
\endlastfoot
$\Omega \subset \mathbb{R}^d$ & Open, bounded domain with Lipschitz boundary $\partial\Omega$, or flat torus $\mathbb{T}^d$. \\
$L^p(\Omega; \mathbb{R}^{d_v})$ & Lebesgue space of measurable functions with norm $\|u\|_{L^p} = (\int_\Omega |u(x)|^p dx)^{1/p}$. \\
$H^s(\Omega) = W^{s,2}(\Omega)$ & Sobolev space of functions whose weak derivatives up to order $s$ reside in $L^2(\Omega)$. \\
$C^k(\overline{\Omega})$ & Space of $k$-times continuously differentiable functions equipped with uniform norm. \\
$X, Y, \mathcal{A}, \mathcal{U}$ & Real or complex Banach or Hilbert spaces. \\
$\mathcal{L}(X, Y)$ & Space of bounded linear operators from $X$ to $Y$ with operator norm $\|T\|_{\mathcal{L}(X,Y)}$. \\
$\langle \cdot, \cdot \rangle_{X^*, X}$ & Duality pairing between Banach space $X$ and its continuous topological dual $X^*$. \\
$\mathcal{F}, \mathcal{F}^{-1}$ & Continuous or discrete Fourier transform and its inverse. The convention is fixed once: $\hat{u}(k) = \int u(x)\, e^{-2\pi i k \cdot x}\,dx$, so $\{e^{2\pi i k \cdot x}\}$ is the $L^2$-orthonormal basis on the torus, and the DFT is unnormalized forward with $1/N$ on the inverse --- PyTorch's default \texttt{norm="backward"} --- unless a passage says otherwise. The aliasing identity of Part~I, the Parseval statements in Appendix~A, and the coefficient magnitudes archived in Appendix~B all turn on this choice. \\
$\mathbb{T}^d \cong [0, 1)^d$ & $d$-dimensional periodic torus identifying opposite boundaries. \\
$D^\alpha u$ & Weak partial derivative for multi-index $\alpha \in \mathbb{N}_0^d$. \\
$\mathcal{M}_N = \{x_j\}_{j=1}^{N}$, $P_N$ & A discretization of $\Omega$ and the restriction operator sending $a$ to its nodal values; \eqref{eq:ch21-grid-layout} and \eqref{eq:ch21-node-layout} are its two array forms. \\
$Z_{k_{\max}}$, $k_{\max}$ & The retained mode set of Part~I, $Z_{k_{\max}} = \{k = (k_1, \dots, k_d) \in \mathbb{Z}^d : |k_j| \le k_{\max, j},\ j = 1, \dots, d\}$ of \eqref{eq:ch05-modeset} --- a box rather than the ball $\{|k| \le k_{\max}\}$, because a product of contiguous index ranges is what a batched transform can slice without gathering scattered entries. It contains $|Z_{k_{\max}}| = \prod_{j=1}^d (2k_{\max, j}+1)$ modes, one complex $d_v \times d_v$ multiplier each, and $k_{\max}$ without a subscript denotes the common cutoff, the case assumed in Theorem~\ref{thm:fno_convergence}. The partial sum $P_{k_{\max}}$ of \eqref{eq:ch05-projection} projects onto the span of $Z_{k_{\max}}$, so it is a spectral projection and not the nodal restriction $P_N$ above, and \eqref{eq:ch05-projection-bound} sums its tail over $k \notin Z_{k_{\max}}$, exactly the modes the box projection discards --- an earlier drafting summed over the ball, whose relation to the box complement runs only one way. Part~II's $Z_m$ of \eqref{eq:ch23-corners} is this same object realized on the half-spectrum, where the two corners carry the cutoffs $m_x, m_y$ in place of the symmetric ranges $|k_j| \le k_{\max, j}$. \\
$\alpha$ & The smoothing order, or regularity gain, of an operator: Theorem~\ref{thm:fno_convergence} assumes $\mathcal{G} : H^s(\mathbb{T}^d; \mathbb{R}^{d_a}) \to H^{s+\alpha}(\mathbb{T}^d; \mathbb{R}^{d_u})$ with $\alpha > 0$ on a compact set $K$; that gain, composed with the compactness, bounds the output regularity $s+\alpha$ that sets the exponent of the truncation term $C_1 k_{\max}^{-(s+\alpha)}$ in \eqref{eq:ch05-error-decomposition} --- the exponent is the regularity of the target \emph{field}, not the gain alone, per \eqref{eq:ch05-projection-bound}. It is a hypothesis about the operator being learned, not a property of the architecture. Not the diffusivity $\alpha = 0.01$ of the periodic heat equation \eqref{eq:appb-pde} in Appendix~\ref{app:b}: there $\alpha$ is a coefficient of the equation, and the mode damping $e^{-\alpha (2\pi k)^2 T}$ of its solution operator \eqref{eq:appb-multiplier} decays faster than any power of $k$, so that number is not the smoothing order of the operator it defines. Nor the multi-index $\alpha \in \mathbb{N}_0^d$ of $D^\alpha u$ above. Three further readings are in use locally: the H\"older exponent of $C^{j,\alpha}$ in the supercritical Sobolev embedding, the miss level $\alpha$ of a conformal interval of coverage $1-\alpha$ in Part~III, and the rollout-horizon multiple $\alpha K$ of Part~II's pushforward training. \\
\midrule
\multicolumn{2}{l}{\textit{Tensor-layout conventions of Part~II}}\\
\addlinespace
$B$ & Batch size; the leading axis of every tensor, and the only one carrying no operator semantics. \\
$C$ & Channel axis of a grid tensor --- the $d_a$ input or $d_u$ output components of Part~I, per node. \\
$C_{\text{in}}, C_{\text{out}}$ & Channel counts at the two ends of one layer; \eqref{eq:ch21-augment} raises $C_{\text{in}}$ to $C_{\text{in}} + d$. \\
$H, W$ & Height and width of a two-dimensional grid; in $d$ dimensions they are the $N_1, \dots, N_d$ of \eqref{eq:ch21-grid-layout}. \\
$d_v$, $L$ & The lifted channel width --- Part~I's $d_v$ carried through the network as \texttt{hidden\_channels} --- and the number of operator layers. \\
$N$, $N_{\text{nodes}}$ & Node count $\prod_j N_j$ of a discretization, and the per-sample node count of a mesh layout. \\
$X \in \mathbb{R}^{B \times C \times N_1 \times \cdots \times N_d}$ & Structured-grid (channel-first) layout; $\prod_j N_j = N$, $C = d_a$ at the input and $C = d_u$ at the output. \\
$X_{\text{nodes}} \in \mathbb{R}^{B \times N_{\text{nodes}} \times C}$ & Point-cloud (mesh) layout, with $X_{\text{coords}} \in \mathbb{R}^{B \times N_{\text{nodes}} \times d}$ carrying the positions $x_j$. \\
$w_j$, $\Delta x$ & Quadrature weights --- cell volumes or Voronoi dual areas --- and the spacing of an equispaced grid, $w_j = \Delta x^d$. \\
$m_x, m_y$ & Retained Fourier modes along $H$ and $W$; written $k_x, k_y$ in Figure~\ref{fig:fno_spectral} and \texttt{modes1}, \texttt{modes2} in the code. \\
$Z_m$, $R^{(1)}, R^{(2)}$ & The retained mode set of \eqref{eq:ch23-corners} and the two complex weight tensors of \eqref{eq:ch23-Rtensors}, indexed by mode, not by grid point. \\
$m$, $P$, $p$, $h$ & DeepONet's sensor count, query-point count, latent width (\texttt{latent\_p}), and hidden width (\texttt{hidden\_dim}). \\
$r$, $\mathcal{N}_r(i)$, $N_{\text{edges}}$ & GNO neighborhood radius, the ball $\mathcal{B}_r(x_i)$ as a node set, and the edge count of \texttt{edge\_index}. \\
$M$, $\mu(x)$, $\sigma(x)$ & Training-ensemble size and the pointwise mean and standard deviation of \eqref{eq:ch21-normalize}; the $\mu$ elsewhere in the book is the input \emph{measure} they estimate. \\
\midrule
\multicolumn{2}{l}{\textit{Physical symbols of Part~III}}\\
\addlinespace
$Q$, $Q^{\star}$ & Query count of a multi-query workload and its break-even value, \eqref{eq:ch31-ledger} and \eqref{eq:ch38-crossover}; Part~I writes the same quantity $M$ in \eqref{eq:ch01-cost-bookkeeping}. \\
$\mathbf{u}$, $\rho$, $\mu$, $\mu_t$ & Velocity and density of the classical solvers Part~III compares against, their molecular viscosity $\mu$ and eddy viscosity $\mu_t$; a subscript $\alpha$ on $\mu_\alpha$ or $\rho_\alpha$ labels a \emph{phase}, not the measure $\mu$ or the mean $\mu(x)$ above. \\
$\mathbf{K}(x)$, $\phi(x)$ & Heterogeneous permeability tensor and porosity of the two-phase Darcy system \eqref{eq:ch34-cons-w}; $\phi$ there is porosity, while $\phi$ in \eqref{eq:ch33-momentum} is latitude. \\
$S_w$, $S_g$, $S_{\mathrm{CO}_2}$ & Phase saturations with $S_w + S_g = 1$; $S_{\mathrm{CO}_2}$ is the reservoir-engineering name for the gas saturation $S_g$, used in Figure~\ref{fig:vol3_sub}. \\
$k_{r\alpha}$, $\lambda_\alpha$, $f_g$, $p_c$ & Relative permeability in $[0,1]$ from the van Genuchten or Brooks--Corey closure, the phase mobility $\lambda_\alpha = k_{r\alpha}/\mu_\alpha$, the fractional flow $f_g$ of \eqref{eq:ch34-mobility}, and the capillary pressure $p_g - p_w$. \\
$\mathbf{u}_\alpha$, $q_\alpha$, $s$ & Darcy phase flux \eqref{eq:ch34-flux}, the well source or sink, and the front speed of the jump condition \eqref{eq:ch34-rh}, across which $[\,\cdot\,]$ denotes the jump. \\
$V_m$, $C_m$, $\chi$, $I_{\text{ion}}$, $I_{\text{stim}}$ & Transmembrane voltage, membrane capacitance per unit area, surface-to-volume ratio, and the ionic and stimulus current densities of the monodomain equation \eqref{eq:monodomain}. \\
$\mathbf{w}$, $k$, $a$, $b$, $\epsilon_0$, $\mu_1$, $\mu_2$ & Recovery variable of \eqref{eq:ch35-recovery} --- distinct from the quadrature weights $w_j$ above --- and the dimensionless Aliev--Panfilov kinetic constants of \eqref{eq:ch35-ion}; the $k$ there is a kinetic parameter, not the mode index of \eqref{eq:ch05-modeset} or the turbulent energy below. \\
$\mathbf{D}(x)$, $\sigma_l$, $\sigma_t$, $\mathbf{f}_0$ & Cardiac diffusion tensor \eqref{eq:ch35-dtensor} built from longitudinal and transverse conductivities along the fiber direction $\mathbf{f}_0$; $\mathcal{D} = \mathbf{D}/(\chi C_m)$ and $D_{\max} = \max_x \rho_{\max}(\mathcal{D}(x))$ set the step \eqref{eq:ch35-cfl}. \\
$h$, $\Delta t$, $T$, $W(h,T)$ & Mesh width, time step, simulated horizon, and the discretization work \eqref{eq:ch35-cost}; $h$ here is a length in micrometers, not the DeepONet hidden width of the block above. \\
$Re$, $U_\infty$, $L$, $\delta$, $y^{+}$, $u_\tau$, $\tau_w$ & Reynolds number \eqref{eq:ch32-reynolds}, freestream speed, body length, boundary-layer thickness, wall-normal distance in viscous units and friction velocity \eqref{eq:ch32-yplus}, and wall shear stress \eqref{eq:ch32-wss}. \\
$C_d$, $C_l$, $F_d$, $F_l$, $A_{\text{ref}}$ & Drag and lift coefficients with their forces and reference area, defined by the surface integral \eqref{eq:ch32-force-coeff}; $k$ and $\omega$ in the closure \eqref{eq:ch32-reynolds-stress} are turbulent kinetic energy and specific dissipation. \\
$\mathbf{u}_h$, $w$, $\theta$, $f$, $\Omega_{\text{rot}}$ & Horizontal and vertical wind, potential temperature, Coriolis parameter $f = 2\Omega_{\text{rot}}\sin\phi$, and planetary rotation rate of \eqref{eq:ch33-momentum}--\eqref{eq:ch33-pottemp}; $\theta$ there is a field, not the network parameter of Part~II. \\
$\mathcal{R}$, $\mathbf{M}$, $\mathcal{A}$, $p(\theta,s)$ & Radon transform, restriction to the $M$ measured view angles, sampled forward operator $\mathcal{A} = \mathbf{M}\mathcal{R}$, and sinogram of \eqref{eq:ch36-radon}, with detector position $s$, detector frequency $\omega$, and moments $m_k$ \eqref{eq:ch36-moment}; this $M$ counts views, not ensemble members. \\
$\mathrm{ACC}$, $\mathrm{RQE}$, $\mathcal{T}$ & Anomaly correlation \eqref{eq:ch33-acc} and relative quantile error \eqref{eq:ch33-rqe}; $\mathcal{T}(w) = \int_\Omega w\,\mathrm{d}x$ the conserved total whose defect $\Delta_{\mathcal{T}}$ \eqref{eq:ch38-defect} measures. Image quality is reported as PSNR, SSIM, and RMSE in Hounsfield units (HU), \eqref{eq:ch36-metrics}. \\
\end{xltabular}
}

% NOTE: the copy of the notation table that once opened Ch. 1 was removed; this
% front-matter table is the single reference for the whole book.
% The Part II tensor-layout block and the Part III physical block extend it with
% the symbols as actually printed in the chapters they name. The two closing
% rows of the Part I block, for the retained mode set and the smoothing order,
% follow eq. (5.5) and Theorem 5.2 of Ch. 5 and the use of those symbols in
% Ch. 2, Ch. 6, Appendix A and Appendix B.
